% Generated draft: requires mathematical and source review.
\documentclass[11pt]{article}
\usepackage[margin=25mm]{geometry}
\usepackage{fontspec}
\setmainfont{Latin Modern Roman}
\usepackage{amsmath,amssymb,mathtools,unicode-math}
\setmathfont{Latin Modern Math}
\usepackage{enumitem,needspace,xurl,hyperref}
\setlength{\parindent}{0pt}
\setlength{\parskip}{5pt}
\setlength{\emergencystretch}{3em}
\newcommand{\N}{\mathbb{N}}
\newcommand{\Z}{\mathbb{Z}}
\newcommand{\Q}{\mathbb{Q}}
\newcommand{\R}{\mathbb{R}}
\newcommand{\C}{\mathbb{C}}
\newcommand{\F}{\mathbb{F}}
\newcommand{\A}{\mathbb{A}}
\newcommand{\PP}{\mathbb{P}}
\newcommand{\E}{\mathbb{E}}
\newcommand{\eps}{\varepsilon}
\newcommand{\dd}{\,\mathrm{d}}
\begin{document}
\section*{UnsolvedMath ID: 20003345}
\textbf{Source problem code:} catalogue entry
\subsection{Definitions and mathematical statement}
The canonical record is a historical-status fragment, not the statement of Conjecture 48: \emph{Conjecture 48 holds when \(\dim D=1\) (see also [38]). The case \(D=A\) was settled in [188]. Recently,} The final sentence is incomplete. Neither Conjecture 48 nor the definitions of \(D\), \(A\), its numbered clauses, or its group hypotheses are present. The attached aid attributes the conjecture to Weiss, the one-dimensional results to Weiss and Dani, and the full maximal real split torus case to Tomanov--Weiss. It cautions that the one-dimensional claim needs an almost-\(\Q\)-simple or suitable factorwise scope: the missing clause (3), read literally for general semisimple products, is reported false. This record cannot specify that clause or certify an unrestricted formulation.\par A separate companion statement supplied by the aid has a precise elementary lattice formulation. Let \(X=\mathrm{SL}_2(\R)/\mathrm{SL}_2(\Z)\) be the space of covolume-one lattices \(L=g\Z^2\subset\R^2\), and set \(a_t=\operatorname{diag}(e^t,e^{-t})\). A ray \(a_tL\), \(t\to+\infty\), is divergent if it eventually leaves each compact subset of \(X\); this is stronger than an unbounded orbit. The aid's claimed criteria are\[a_tL\text{ is divergent as }t\to+\infty\ \Longleftrightarrow\ L\cap(\{0\}\times\R)\ne\{0\},\qquad a_tL\text{ is divergent as }t\to-\infty\ \Longleftrightarrow\ L\cap(\R\times\{0\})\ne\{0\}.\] Consequently the full orbit map \(\R\to X\), \(t\mapsto a_tL\), is proper if and only if both a nonzero horizontal vector and a nonzero vertical vector belong to \(L\). Equivalently, \(L\) contains a finite-index rectangular sublattice \(L_0=\Z(x,0)+\Z(0,y)\) for some nonzero real \(x,y\). With the covolume-one convention, \([L:L_0]=|xy|\) is a positive integer. For \(L=g\Z^2\), primitive integral vectors \(m_+,m_-\in\Z^2\) can be chosen with \(gm_+\) vertical and \(gm_-\) horizontal. They are independent rational vectors in the standard representation and satisfy \(a_tgm_+\to0\) as \(t\to+\infty\) and \(a_tgm_-\to0\) as \(t\to-\infty\). These are the two rational witnesses for the aid's description of obvious divergence. The forward-ray criterion says that a single primitive lattice line suffices throughout a divergent ray, not merely that a different short vector can be found at each time. These companion assertions do not supply the missing general Conjecture 48.
\subsection{Short English statement}
This item records two historically settled cases of an omitted conjecture. The attached aid gives a separate two-dimensional criterion: a diagonal ray diverges precisely when the lattice has a vector in its contracting coordinate direction, and the full orbit map is proper precisely when both coordinate directions occur.
\subsection{Sources}
\begin{itemize}\item Local UnsolvedMath registry, permanent ID 20003345, code catalogue entry; canonical \texttt{aim-other-notes.json}, \texttt{notes[151]}, conjecture tag.\item American Institute of Mathematics Problem Lists, Other domain, \emph{Emerging applications of measure rigidity}, source item 48, historical paragraph concerning Conjecture 48; no section supplied. \url{https://aimath.org/WWN/measrigid/measrigid.pdf}.\item Source references [38] and [188]; complete bibliographic entries are absent. The attached aid attributes the conjecture to Weiss, one-dimensional results to Weiss and Dani, and the maximal real split torus case to Tomanov--Weiss.\item Attached research attempt 1, \emph{The one-parameter and full-torus status cases of Conjecture 48}, recorded 2026-08-14: \texttt{context\_only}, problem-at-run \texttt{not\_a\_problem}, rigor \texttt{proved}, novelty confidence low; registry label \texttt{open}, classification \texttt{OPEN-TRIAGE}. The machine-generated aid is not peer reviewed and requires independent verification before citation.\end{itemize}
\end{document}
